\section{Elementary primitives: the finite reduction}
\label{sec:rational-primitives}

The target theorem states that a rational function on an interval whose
denominator is uniformly apart from zero has an elementary primitive.
Algebraic constants are allowed. The general analytic theorem remains open;
the checked results here are finite formal-derivative identities and a
domain-preserving circle substitution.

\begin{theorem}[Supplied partial fractions]
\label{thm:rational-primitive-formula}
\lean{ComputableAnalysis.RationalPrimitiveFormula.Decomposition.primitive_correct}
\leanok
Given a certified decomposition with rational coefficients into a polynomial,
linear poles and positive quadratic poles, the finite primitive algorithm
returns a formula whose formal derivative is the input rational function at
every permitted rational input.
\end{theorem}
\begin{proof}\leanok
For \(Q=(x-a)^2+b\), \(b>0\), repeated quadratic poles reduce by
\[
 I_1=\frac1{\sqrt b}\arctan\frac{x-a}{\sqrt b},\qquad
 I_{m+1}=\frac{x-a}{2mbQ^m}+\frac{2m-1}{2mb}I_m.
\]
Linear numerators reduce to logarithms and reciprocal powers of \(Q\).
Linear poles give logarithms or reciprocal powers of \(x-a\); polynomial
coefficients integrate termwise. Induction checks the finite identities.
This is a statement about formal differentiation, not a derivative
certificate for a represented elementary evaluator.
\end{proof}

\begin{theorem}[Rational circle pullback]
\label{thm:rational-circle-pullback}
\lean{ComputableAnalysis.TrigonometricRationalization.Expr.pullback_correct,
ComputableAnalysis.TrigonometricRationalization.Expr.pullback_defined_iff}
\leanok
Every rational expression in two circle coordinates pulls back to a rational
function under
\[
 s=\frac{2t}{1+t^2},\qquad c=\frac{1-t^2}{1+t^2},
 \qquad J=\frac2{1+t^2}.
\]
The computed rational function equals the expression times \(J\), and is
defined exactly when the original expression at those coordinates is defined.
\end{theorem}
\begin{proof}\leanok
Structural induction on rational expressions and domain-preserving polynomial
normalization. The reciprocal of \(p/q\) is stored as \(q^2/(pq)\), so the
original poles are not filled by cancellation.
\end{proof}

Two antipodal charts cover the rational circle. General algebraic partial
fractions, interval-valued elementary realizations, analytic differentiation,
and transport through the represented angle charts are still needed for the
full rational-primitive theorem and its trigonometric corollary.
